\documentclass[11pt,letterpaper]{article}
\usepackage[margin=0.9in,headheight=14pt]{geometry}
\usepackage{amsmath,amssymb}
\usepackage{enumitem}
\usepackage{fancyhdr}
\usepackage{lmodern}
\usepackage[T1]{fontenc}
\usepackage{graphicx}
\usepackage{xcolor}

\pagestyle{fancy}
\fancyhf{}
\lhead{SYDE 556/750 \quad Practice Test 3B: Nengo \& Simple Dynamics}
\cfoot{\thepage}
\renewcommand{\headrulewidth}{0.4pt}

\setlength{\parindent}{0pt}
\setlength{\parskip}{5pt}

\newcommand{\trc}{\tau_{RC}}
\newcommand{\tref}{\tau_{\mathrm{ref}}}
\newcommand{\question}[2]{\par\vspace{10pt}\textbf{Question #1.}\hfill\textbf{[#2 marks]}\par}
\newcommand{\answerbox}[1]{\par\noindent\framebox[\linewidth]{\rule{0pt}{#1}}\par}

\setlist[itemize]{itemsep=1pt}

\begin{document}

\begin{center}
\IfFileExists{../lectures/assets/uwlogo.pdf}{\includegraphics[height=1.1cm]{../lectures/assets/uwlogo.pdf}\\[10pt]}{}
SYDE 556/750 -- Simulating Neurobiological Systems\\[2pt]
{\Large\bfseries Practice Test 3B: Nengo \& Simple Dynamics}
\end{center}

\vspace{6pt}
Student ID: \rule{4.5cm}{0.4pt} \hfill Time: 30 minutes \hfill Total: 20 marks (plus 2 bonus)

\begin{itemize}
\item \textbf{Closed book.} No notes, textbook, calculator, phone, or other devices.
\item The marks for each question are shown in brackets on the right.
\item Write your answers in the space provided. Show your work --- partial credit is given for correct reasoning.
\item Unless stated otherwise, every synapse is the first-order low-pass $h(t) = \frac1\tau e^{-t/\tau}$.
\end{itemize}
\hrule

%---------------------------------------------------------------------
\question{1}{4}
(a) A Nengo ensemble represents $\mathbf x \in \mathbb R^2$ with \texttt{radius=2}. One of its LIF neurons has encoder $e = (0, 1)$, a (normalized) intercept of $-0.5$, and a maximum rate of $150$ Hz. In the left box, sketch this neuron's firing rate along the line $\mathbf x = (0, s)$, and along the line $\mathbf x = (s, 0)$, for $s \in [-2, 2]$. Label where it starts firing.

(b) The same kind of ensemble (now 1-D, radius 1) is driven by $x(t) = \sin(2\pi \cdot 2t)$. Its decoded output is recorded with three probes: \texttt{synapse=0.001}, \texttt{synapse=0.01} and \texttt{synapse=0.1}. In the right box, sketch what each probe records over one second. Does changing the probe synapse change what the ensemble represents?

\noindent\framebox[0.5\linewidth][l]{\rule{0pt}{5.2cm}(a)}\framebox[0.5\linewidth][l]{\rule{0pt}{5.2cm}(b)}

\vspace{2.2cm}

\newpage
%---------------------------------------------------------------------
\question{2}{4}
A population representing $x$ receives an input $u$ through a connection computing $g(u)$, and has a recurrent connection computing $f(x)$. Both connections use a synapse with time constant $\tau$. The represented value obeys
\[
\tau\,\frac{dx}{dt} = g(u) + f(x) - x.
\]
(a) Explain what each term represents, and where the $-x$ comes from. What would $x$ do if the recurrent connection were removed?
\vspace{4.2cm}

(b) Using $\tau = 50$ ms, find the recurrent and input transforms so that the network implements the low-pass filter
\[
\frac{dx}{dt} = -\frac{1}{0.5}\,x + \frac{1}{0.5}\,u.
\]
Draw the network (populations, connections, transforms). What time constant would you see in the response to a step in $u$?
\vspace{5cm}

\newpage
%---------------------------------------------------------------------
\question{3}{4}
A population with radius $1$ and synaptic time constant $\tau$ has input connection $g(u) = u$ and recurrent connection $f(x) = x - x^2$. The input is held constant.

(a) Write the resulting dynamics $\tau\dot x = \varphi(x)$. For $u = 0.25$, find all equilibria and state whether each is stable or unstable.
\vspace{3.4cm}

(b) Sketch $\varphi(x)$ for $u = 0.25$ with arrows on the $x$-axis. Describe what $x$ does starting from $x(0) = 0$ and from $x(0) = -0.8$. What changes if the input is lowered to $u = -0.1$?
\answerbox{5.4cm}

\newpage
%---------------------------------------------------------------------
\question{4}{4}
You want a 2-D population to oscillate at $1$ Hz:
\[
\dot{\mathbf x} = \begin{bmatrix} 0 & -\omega \\ \omega & 0\end{bmatrix}\mathbf x, \qquad \omega = 2\pi \text{ rad/s}.
\]
(a) The recurrent synapse is $\tau = 0.1$ s. What recurrent transform $A'$ do you use? (Leave $\pi$ in your answer.)
\vspace{2.4cm}

(b) A brief input kick leaves the state at $\mathbf x(0) = (0, 0.8)$. In the box, sketch \emph{all of the following}:
\begin{itemize}
\item $x_1(t)$ and $x_2(t)$ on the same time axis for $0 \le t \le 2$ s;
\item the trajectory in the $(x_1, x_2)$ plane, with an arrow showing its direction.
\end{itemize}
\answerbox{5cm}

(c) Why does the network need the kick? In a real Nengo simulation, what sets the amplitude it ends up oscillating at?
\vspace{2.4cm}

\newpage
%---------------------------------------------------------------------
\question{5}{4}
An integrator $\dot x = u$ is built with recurrent transform $1$ and recurrent synapse $\tau = 0.1$ s. It starts at $x(0) = 0$ and receives a ramp input:
\[
u(t) = \begin{cases} 4t, & 0 \le t < 0.5 \text{ s},\\ 0, & t \ge 0.5 \text{ s}.\end{cases}
\]
(a) Find the ideal $x(t)$ for all $t \ge 0$ and sketch it.
\vspace{2.2cm}
\answerbox{3.6cm}

(b) The input connection uses transform $\tau$ but a \emph{faster} synapse, $\tau_{\mathrm{in}} = 10$ ms. Sketch how the decoded $x(t)$ differs from your answer to (a). Roughly how large is the difference just before $t = 0.5$ s, and what happens after the ramp stops?
\answerbox{3.6cm}
\vspace{1.6cm}

\newpage
%---------------------------------------------------------------------
\textbf{Bonus Question 1.}\hfill\textbf{[2 bonus marks]}\par
You build the integrator of Question 5 (with matched synapses) and stop the input at $t = 0.5$ s. Over the next few seconds, the decoded value creeps away from $0.5$. In plain language, why does a neural integrator drift? With spiking neurons the drift looks jittery and partly random, but with the same neurons in rate mode (\texttt{LIFRate}) it is smooth and predictable. Why?
\vspace{5cm}

\end{document}
